Studierende planen ihr Studium anhand von Curricula, die Voraussetzungen,
Wahlmöglichkeiten und institutionelle Regeln festlegen. Die Werkzeuge, die diese
Planung unterstützen, sind zweigeteilt. Graphbasierte Systeme stellen die
Struktur eines Curriculums dar, lassen jedoch den zeitlichen Ablauf außen vor;
tabellenbasierte Planungswerkzeuge verteilen Lehrveranstaltungen auf Semester
und prüfen das Ergebnis gegen die Studienordnung, lassen aber die Struktur
implizit. Wo beides in einem Werkzeug zusammengeführt wurde, ist die
Abhängigkeitsinformation in den Zeitplan selbst eingebettet und damit an eine
chronologische Anordnung gebunden.

Diese Arbeit entwirft, implementiert und evaluiert eine hybride
Graph-Tabellen-Oberfläche, in der ein navigierbarer Curriculumsgraph und eine
semesterorientierte Tabelle auf einen gemeinsamen Plan schreiben. Die
Anforderungen stehen am Anfang: Eine formative Studie mit sechs Studierenden,
durchgeführt anhand eines interaktiven, aber nicht funktionsfähigen Prototyps,
ergibt eine gereihte Spezifikation von 13 Features mit 55 nummerierten
Anforderungen. Das Werkzeug wird anschließend mit elf Studierenden evaluiert,
die jeweils einmal ohne und einmal mit der Graphansicht planen und dabei über
eine bildweise Aufzeichnung der Oberfläche, ein Think-aloud-Transkript und
Fragebögen nach den Aufgaben beobachtet werden.

Eine explorative Evaluierung mit elf Studierenden zeigte eine durchgängige, von
einer Checkliste getriebene Planungsschleife über alle Erfahrungsstufen hinweg.
Der Graph trat bei der Suche nach Kandidaten an die Stelle des
Lehrveranstaltungskatalogs, während Lehrveranstaltungen weiterhin in der Tabelle
festgelegt wurden. Dieser geteilte Arbeitsablauf machte eine Parking Stage als
Puffer für noch nicht festgelegte Lehrveranstaltungen zwischen den beiden
Ansichten unverzichtbar. Da das Curriculum zudem nur wenige strikte
Abhängigkeiten festlegt, verschob sich der wesentliche Nutzen des Graphen vom
Abbilden von Reihenfolgen hin zur Orientierung und zum gefilterten Erkunden. Die
berichtete Zufriedenheit war insgesamt hoch, wobei die schwächste Skala
(Perspicuity) genau mit den Terminologieschwierigkeiten übereinstimmt, die
während der Aufgaben beobachtet wurden.
